\section{Introduction}\label{sec:introduction}

For finite simple undirected graphs \(H\) and \(G\) with
\(|V(G)|\ge1\), write
\[
t(H,G)=\frac{\Hom(H,G)}{|V(G)|^{|V(H)|}},
\qquad
p(G)=\frac{2|E(G)|}{|V(G)|^2},
\]
where \(\Hom(H,G)\) counts every edge-preserving map, including
noninjective maps. Sidorenko's conjecture asserts that
\begin{equation}\label{eq:sidorenko}
t(H,G)\ge p(G)^{|E(H)|}
\end{equation}
whenever \(H\) is bipartite and has at least one edge.
We disprove this conjecture.

\begin{theorem}\label{thm:main}
Let \(H\) be the incidence graph of the \(22\) triples on \(13\) points
in Table~\ref{tab:complex}. This bipartite graph has \(35\) vertices
and \(66\) edges. There is a finite simple undirected graph \(G\)
with at least one edge such that
\[
t(H,G)<p(G)^{66}.
\]
\end{theorem}

We first construct a strict density gap for a finite nonnegative
kernel and then retain it when sampling a finite simple host. This
route specifies \(H\) explicitly and proves existence of \(G\), while
the final density continues to count all homomorphisms.

A finite kernel is a nonnegative function \(W(x,y)\) on two finite
probability spaces. Its mean is the analogue of edge density, and its
\(H\)-density is the expectation of the product of \(W\) over the
incidences of \(H\), with independent samples from the first space at
the point vertices and from the second space at the face vertices.
For a symmetric kernel taking values in \([0,1]\), a simple graph
can be sampled by choosing a label for each host vertex and using the
kernel value as the probability of each edge.

The same pattern has further consequences with distinct density
scopes. Tensor powers make \(t(H,F)/p(F)^{66}\) arbitrarily small,
with the host density allowed to vary. Balanced blow-ups retain a
fixed deficit in injective labeled copies at one positive limiting
edge density. At one fixed \(p\in(0,1)\), a nonconstant symmetric
\([0,1]\)-valued kernel has mean \(p\) and \(H\)-density \(p^{66}\);
sampling gives finite graphs whose edge and \(H\)-densities converge
to these values while their four-cycle density has a limit strictly
greater than \(p^4\).
Corollaries~\ref{cor:amplification} and~\ref{cor:forcing} state these
consequences precisely below.

\subsection*{The negative correlation}

The triple system has \(13\) points and \(22\) faces. Each of its
\(33\) point pairs belongs to exactly two faces. This double
occurrence converts a local sign constraint into a correlation
involving four incidences.

To see the algebra, attach a sign \(\xi_{ij}\) to each chosen active
incidence and an independent uniform sign \(\eta_e\) to each point
pair. The finite-field lemma compares the averaged matrix product with
a sign model having a factor
\(1+c_0\eta_e\xi_{ij}\xi_{kj}\) whenever both incidences of
\(e=\{i,k\}\) on a face \(j\) are active; here
\(c_0\in\{-1,1\}\). If the two faces on \(e\) are \(j,l\), then
\[
 \E_{\eta_e}
 (1+c_0\eta_e\xi_{ij}\xi_{kj})
 (1+c_0\eta_e\xi_{il}\xi_{kl})
 =
 1+\xi_{ij}\xi_{kj}\xi_{il}\xi_{kl}.
\]
A factor used on only one face averages to \(1\). Thus the model
retains the parity of the four incidence signs exactly when both
face occurrences are selected. This is the finite correlation that
the kernel construction turns negative.

The matrix model uses a large fixed even dimension \(D=2r\).
For a symmetric matrix difference over \(\F_q\), with \(q\) an odd
prime, prescribe rank \(r\)
and the determinant sign of the induced nondegenerate quotient form.
Half rank is the key geometric constraint: when the difference
between two point matrices is nonsingular, the two rank-\(r\) images
through a common center are complementary, which relates their
determinant signs. On configurations with nonsingular point
differences, the normalized center integrals approximate the sign
model in averaged \(L^1\), with a uniform bound on the local factors.
A joint estimate shows that the determinant signs on the \(33\) point
pairs are asymptotically independent, supplying the independent
\(\eta_e\) in the limiting sign model.

The kernel uses one independent matrix coordinate for each incidence.
A type names the point or face that a sampled vertex is intended to
represent. At an incident type pair, two choices of active coordinate
set have opposite sign biases. They cancel every nonconstant first
sign moment while retaining a signed overlap moment. Expanding the
parity model assigns a set of selected coordinates to each point
pair. The chosen supports permit exactly one collection of nonempty
pair labels at the intended type assignment, and its product of
overlap moments is negative. Scores obtained from color refinement
make this assignment dominate all other part-preserving type maps,
including maps with repetitions.

The limiting density in one orientation now has a negative correction.
Adding an inactive point type makes that correction occur with a factor
\(\varepsilon^{13}\), while the uncontrolled correction in the
opposite orientation has a factor \(\varepsilon^{22}\). Choosing
\(\varepsilon>0\) sufficiently small and then taking \(q\) sufficiently
large makes the product of the two finite oriented densities less than
one. A crossed product of the rectangular kernel
with its transpose then gives a symmetric kernel with the same strict
deficit.

The remaining issue is the contribution of singular point
differences. Their probability is \(O(q^{-1})\), but the normalized
rank indicators can grow as a power of \(q\), so the weighted
contribution requires a separate bound. The graph
\(\{(v,Xv):v\in\F_q^D\}\) of a symmetric matrix is a Lagrangian
subspace of \(\F_q^D\oplus\F_q^D\) for
\[
 \omega((v,w),(v',w'))=v^{\mathsf T}w'-w^{\mathsf T}v'.
\]
The intersection of two such graphs projects isomorphically onto
the kernel of their difference. Rank defects thereby become
intersection dimensions.

We compare the probability of each point-intersection profile with
the gain from conditioning on its intersections with a face center.
Three partitions of the face set give a global profile bound. If this
first bound does not give decay, a quadratic estimate forces any large
pair intersection to propagate through the connected complex until
the extra separation of one pair makes the profile too costly. The
remaining bounded dimensions reduce, after choosing \(D\), to a
short list of integer equality cases. A final span count gives strict
decay for those cases. The dimension and all finite construction
parameters are fixed before \(q\) tends to infinity; constants in
the \(q\)-estimates may depend on that fixed dimension.

Figure~\ref{fig:proof-overview} summarizes these dependencies.

\begin{figure}[htbp]
\centering
\begin{tikzpicture}[
  font=\small,
  proofbox/.style={draw=linkblue!70,rounded corners=2pt,
    fill=linkblue!4,align=center,inner sep=5pt,text width=3.4cm},
  proofstep/.style={-{Latex[length=2mm]},thick,draw=linkblue!80}]
  \node[proofbox,text width=6.0cm] (complex) at (0,0)
    {Fixed triple system\\Section~\ref{sec:complex}};
  \node[proofbox,text width=3.8cm] (activation) at (-3.7,-1.5)
    {Activation supports\\and type scores\\Section~\ref{sec:kernel}};
  \node[proofbox] (transverse) at (0.9,-1.5)
    {Transverse contribution\\Section~\ref{sec:transverse}};
  \node[proofbox] (tail) at (4.9,-1.5)
    {Weighted singular\\bound\\Sections~\ref{sec:orbits}--\ref{sec:tail}};
  \node[proofbox,text width=5.8cm] (rank) at (2.9,-2.9)
    {Rank-layer limit\\Lemma~\ref{lem:rank-layer}};
  \node[proofbox,text width=6.5cm] (negative) at (0,-4.3)
    {Negative averaged coefficient\\in one limiting oriented density};
  \node[proofbox,text width=6.5cm] (kernel) at (0,-5.6)
    {Dilution and crossed product\\symmetric kernel with a strict deficit};
  \node[proofbox,text width=6.5cm] (host) at (0,-6.9)
    {Sampling a finite simple host $G$\\Lemma~\ref{lem:finite-host}};
  \coordinate (estimatefork) at (2.9,-0.7);
  \draw[proofstep] (complex.south west) -- (activation.north);
  \draw[proofstep] (complex.south east) -- (estimatefork) -- (transverse.north);
  \draw[proofstep] (estimatefork) -- (tail.north);
  \draw[proofstep] (transverse.south) -- (rank.north west);
  \draw[proofstep] (tail.south) -- (rank.north east);
  \draw[proofstep] (activation.south) -- (negative.north west);
  \draw[proofstep] (rank.south) -- (negative.north east);
  \draw[proofstep] (negative) -- (kernel);
  \draw[proofstep] (kernel) -- (host);
\end{tikzpicture}
\caption{Logical dependencies of the counterexample. The fixed triple
system supplies the combinatorial properties used in both branches.
The rank-layer limit realizes the averaged coefficient made negative
by the finite activation construction. Section~\ref{sec:kernel}
constructs the kernel and transfers its strict deficit to a finite
simple host assuming Lemma~\ref{lem:rank-layer};
Sections~\ref{sec:transverse}--\ref{sec:tail} then prove that input.}
\label{fig:proof-overview}
\end{figure}


\subsection*{History and relation to positive results}

The conjecture grew out of questions about supersaturation: how many
copies of a fixed bipartite graph must occur once the host has many
edges? In a related study of graphcopy functions, Erd\H{o}s,
Lov\'asz, and Spencer distinguished normalized homomorphism and
injective-copy counts and proved their limiting agreement under
increasing host orders~\cite{erdoslovaszspencer1979}.
Simonovits's discussion of degenerate extremal problems placed
supersaturation lower bounds in a general program~\cite{simonovits1984}.
Sidorenko formulated the associated inequalities for functionals of
nonnegative kernels~\cite{sidorenko1991,sidorenko1993}.  In the
homomorphism formulation, the proposed lower bound is the density in
a constant kernel of the same mean.  The insistence on all
homomorphisms is part of this formulation; repeated images of vertices
are allowed.  Theorem~\ref{thm:main} concerns this ordinary graph
inequality.

Trees, even cycles, and complete bipartite graphs are basic positive
examples~\cite{sidorenko1993}. Several methods established broader
classes. Conlon,
Fox, and Sudakov proved the conjecture when one vertex is adjacent to
the entire opposite part, and obtained more general comparisons with
complete bipartite graphs~\cite{conlonfoxsudakov2010}.  A different
approach studies norms derived from homomorphism densities.
Hatami developed the graph-norm approach, recording Szegedy's observation
that weakly norming graphs satisfy Sidorenko's inequality, and proved
stronger subgraph inequalities and the hypercube case~\cite{hatami2010}.
Conlon and Lee constructed weakly norming graphs from finite reflection
groups~\cite{conlonlee2017}.  These results provide stronger inequalities
than the comparison with a single edge and explain why many highly
structured bipartite graphs satisfy the conjecture.

Other proofs assemble a graph from pieces while keeping control of
the associated homomorphism distribution.  Li and Szegedy's logarithmic
calculus, first circulated in 2011, gives gluing results and a class of
reflection trees~\cite{liszegedy2025}.  Szegedy developed an information
theoretic formulation using entropy and conditional
gluing~\cite{szegedy2015entropy}.  Kim, Lee, and Lee proved the conjecture
for tree-arrangeable graphs and for Cartesian products of a Sidorenko
graph with a tree~\cite{kimleelee2016}.  Conlon, Kim, Lee, and Lee
introduced strong tree decompositions and established further
subdivision results~\cite{conlonkimleelee2018}.  More recently, Im,
Li, and Liu treated subdivisions and substitutions by even theta
graphs under a local-density counting hypothesis for the base or suitable
global path-multiplicity conditions~\cite{imliliu2026}.
The gluing and substitution statements impose different structural
conditions on their constituent graphs.

Positive results also depend on how often neighborhoods of a given
size occur. Conlon and Lee proved a divisibility criterion and a
degree-count criterion when one part is regular, and showed that every
bipartite graph has a sufficiently large one-sided blow-up satisfying
the conjecture~\cite{conlonlee2021}. Coregliano subsequently replaced
the divisibility hypothesis with weaker degree-count bounds and
extended the reflection approach through left-cut-percolation and
induced-Sidorenko bigraphs~\cite{coregliano2024}.  A positive result
for a blow-up does not imply the corresponding result for the
original graph.

There are useful reductions of the host space as well. Coregliano and
Razborov proved that, for a fixed pattern, a lower bound by one fixed
positive multiple of the Sidorenko expression on every biregular
bigraphon implies the exact Sidorenko inequality on every bigraphon.
Here biregular means that the
kernel integrals in either variable are constant. They also proved a
reflective tree-decomposition theorem~\cite{coreglianorazborov2021}.
In another direction, Zhao gave a conditional reduction through
conjugacy-class averaging on finite groups and proved an inequality
for one-subdivisions on the resulting conjugacy-averaged bipartite
Cayley kernels~\cite[Theorems~1.3--1.4]{zhao2026conjugacy}.  The averaging monotonicity in the
general reduction is a hypothesis.

The incidence graph used here has parts of sizes $13$ and $22$,
with degrees $4,5,6$ on the first part and degree $3$ on the second.
It therefore has no vertex adjacent to the whole opposite part.
Since it is connected and its first part is not regular, it is not
weakly norming: Hatami's necessary condition requires regularity on
each part~\cite{hatami2010}. Coregliano's sufficient degree-count
hypothesis is not met in either orientation.  On the degree-$3$ side the relevant comparison is
$22<\binom{13}{3}=286$; in the reverse orientation there are only
$3,6,4$ vertices of degrees $4,5,6$, respectively, below the
corresponding thresholds
$\binom{22}{4},\binom{22}{5},\binom{22}{6}$~\cite[Theorem~3.6]{coregliano2024}.
Its minimum degree is $3$, so it is not directly a graph obtained by
the even-path substitutions above, which introduce vertices of
degree $2$.  These observations locate the example relative to
specific positive theorems; the negative inequality is established by
the construction and estimates below.

Finally, the conclusion is compatible with local versions of the
conjecture.  For a bipartite graph with $m$ edges, Lov\'asz proved the
inequality for a mean-one symmetric kernel $W$ satisfying
$\lVert W-1\rVert_\infty\leq 1/(4m)$~\cite[Theorem~5.1]{lovasz2011local}.
He also proved a cut-norm
version under the additional bound $0\leq W\leq2$, with
$\lVert W-1\rVert_\square\leq2^{-8m-2}$~\cite[Theorem~4.1]{lovasz2011local}.
Here the kernels are on $[0,1]^2$, with Lebesgue measure, and
\[
 \lVert F\rVert_\square
 =\sup_{S,T\subseteq[0,1]}
   \left|\int_{S\times T}F(x,y)\,dx\,dy\right|,
\]
where $S,T$ are measurable. Thus the constant kernel has a neighborhood
on which the inequality holds. A finite probability kernel can be
represented on $[0,1]^2$ by a step function whose interval lengths
are its atom probabilities. The mean-one normalization of our final
symmetric kernel, having a strict deficit, cannot satisfy either set
of local hypotheses above. Scaling the kernel into $[0,1]$ preserves
the normalized density comparison.

The construction also uses established geometric and algebraic
frameworks in specific roles. Conditions on the rank and
discriminant sign of a symmetric matrix difference are the relations
in the association scheme of symmetric matrices studied by Huo and
Wan~\cite{huowan1993}; Schmidt records this framework and the exact
counts by rank and sign~\cite[Section~2]{schmidt2015}. Symplectic geometry
provides the elementary Lagrangian and pair-intersection counting
formulas~\cite{taylor1992}.
The joint-sign argument uses the standard quadratic Gauss-sum rank
bound~\cite[Lemma~3.3]{gowerswolf2011}, and the type scores adapt
color refinement~\cite{weisfeilerleman1968,mckay1981}. We rederive
the estimates used here and combine the signed rank-layer correlations
with bounds for their weighted contribution on singular configurations.

\subsection*{Amplification and injective copies}

For a finite simple graph \(F\) of order \(n\ge35\), let
\(\operatorname{Inj}(H,F)\)
count injective edge-preserving maps, with no condition on nonedges,
and write \((n)_{35}=n(n-1)\cdots(n-34)\). At ordinary edge density
\(d=|E(F)|/\binom n2\), the labeled random-graph benchmark is
\((n)_{35}d^{66}\), the expected count in \(G(n,d)\).

\begin{corollary}[Amplification and injective supersaturation]
\label{cor:amplification}
For the same graph \(H\) as in Theorem~\ref{thm:main}, the following
hold.
\begin{enumerate}[label=(\roman*)]
\item For every \(c>0\), there is a finite simple undirected graph
\(F\) with \(p(F)>0\) such that
\[
t(H,F)<c\,p(F)^{66}.
\]
\item There are \(p_*\in(0,1)\), \(\delta_*>0\), and finite simple
undirected graphs \(F_s\) of orders \(n_s\to\infty\) for which
\(d_s=|E(F_s)|/\binom{n_s}{2}\to p_*\) and, for all sufficiently
large \(s\),
\[
\operatorname{Inj}(H,F_s)
\le(1-\delta_*)(n_s)_{35}d_s^{66}.
\]
\end{enumerate}
\end{corollary}

The proof in Section~\ref{sec:consequences} uses tensor powers and
balanced blow-ups of a single host from Theorem~\ref{thm:main}.

Part (ii) contradicts the injective Erd\H{o}s--Simonovits
random-graph benchmark described in
\cite[Section~1]{conlonfoxsudakov2010}: the sequence has one positive
limiting density, so it eventually satisfies \(d_s>n_s^{-\gamma}\)
for every fixed \(\gamma>0\), while retaining a fixed deficit.

\subsection*{A consequence for the forcing conjecture}

Fix \(0<p<1\). A graph \(F\) with \(m\) edges is \emph{\(p\)-forcing}
if every sequence of finite simple graphs \(G_n\) with
\(|V(G_n)|\to\infty\), \(p(G_n)\to p\), and
\(t(F,G_n)\to p^m\) is quasirandom of density \(p\): for every fixed
finite simple graph \(J\), one has
\(t(J,G_n)\to p^{|E(J)|}\). It is \emph{forcing} if it is
\(p\)-forcing for every \(p\in(0,1)\). The classical edge-and-four-cycle
characterization of quasirandomness is due to Chung, Graham, and
Wilson~\cite{chunggrahamwilson1989}. Skokan and Thoma studied which
bipartite patterns can replace the four-cycle~\cite{skokanthoma2004}.
Using this asymptotic formulation and following their question,
Conlon, Fox, and Sudakov formulate the forcing conjecture: a graph is forcing if and
only if it is bipartite and contains a cycle
\cite[Section~1, equation~(2) and Conjecture~2]{conlonfoxsudakov2010}.

For a symmetric \([0,1]\)-valued kernel \(W\) on a probability space,
write
\[
t(F,W)=\E\prod_{uv\in E(F)}W(X_u,X_v),
\]
where the \(X_v\) are independent samples. In particular,
\(t(K_2,W)=\E W\). We identify each finite probability kernel below
with its step graphon on \([0,1]\).

\begin{corollary}[Failure of the forcing conjecture]
\label{cor:forcing}
For the same connected bipartite graph \(H\) of
Theorem~\ref{thm:main}, there are \(p\in(0,1)\) and a nonconstant
graphon \(W\) such that
\[
t(K_2,W)=p,\qquad t(H,W)=p^{66}.
\]
Consequently, \(H\) is not \(p\)-forcing at this density and hence is
not forcing, despite being bipartite and containing a cycle.
\end{corollary}

Section~\ref{sec:consequences} interpolates between finite kernels of
the same mean to obtain the exact graphon equality, then uses independent
sampling to obtain a sequence of finite graphs with the stated limiting
densities.
Conlon, Fox, and Sudakov note that the forcing conjecture is stronger
than Sidorenko's conjecture~\cite[Section~1]{conlonfoxsudakov2010};
the interpolation gives a concrete witness for this graph at one density.

\subsection*{Organization and conventions}

Section~\ref{sec:complex} gives the triple system and verifies its
finite properties. Section~\ref{sec:kernel} states
Lemma~\ref{lem:rank-layer} and uses it to prove
Theorem~\ref{thm:main}. The next four sections prove that lemma:
Section~\ref{sec:transverse} treats transverse matrices;
Section~\ref{sec:orbits} gives the Lagrangian orbit counts;
Section~\ref{sec:planted} bounds the local gains; and
Section~\ref{sec:tail} proves strict decay of the singular
contribution. Section~\ref{sec:consequences} then proves the amplification,
injective-copy and forcing consequences stated above.

All probability spaces used in our construction are finite. Expectations
refer to the indicated independent samples, except where a
conditional law is specified. Empty products equal \(1\).
The field size \(q\) tends to infinity through odd primes. Every
such limit is taken after fixing the dimension and all finite
construction parameters. Dependence on \(D\) in the constants of
these \(q\)-asymptotic estimates is allowed; bounds used to choose
\(D\) will explicitly be uniform in \(D\).
